


Generated by \texttt{python -m crashguard.ar\_verdicts}. SUPPORTED needs the preregistered criterion on the held-out data and a Holm-adjusted p-value below 0.05 (family S: 6 hypotheses; family R: 7). Two clean counts for the paper: \textbf{family S (simulated hardware-format validation): 2 supported, 4 not supported of 6. Family R (architecture and protocol candidates): 5 supported, 2 not supported of 7.} These two families are not mixed with the earlier 67 hypotheses of \texttt{PROJECT\_STATUS.md} or with the V5 family of 18.

Everything in Part S is SIMULATED\_HW\_FORMAT: code run on simulated hardware-format logs generated from stated sensor and link models; no real hardware measurements are reported. A candidate enters the main paper only if all its hypotheses are supported; the others go to the supplement, and the failed hypotheses to \texttt{docs/03\_results/NEGATIVES\_REGISTER.md}.

{\scriptsize\raggedright
\begin{longtable}{p{0.086\linewidth}p{0.086\linewidth}p{0.086\linewidth}p{0.086\linewidth}p{0.086\linewidth}p{0.086\linewidth}p{0.086\linewidth}p{0.086\linewidth}p{0.086\linewidth}p{0.086\linewidth}p{0.086\linewidth}}
\toprule
Id & Statement & Label & Criterion met & Raw p & Holm p & Verdict & Effect size and interval & Nearest prior work (Part L) & Allowed wording & Placement \\
\midrule
\endhead
H-S1 & The impact-side rule keeps side accuracy at or above 0.90 in every sensor variant of the simulated logs. & SIMULATED\_HW\_FORMAT & no & 1 & 1 & \textbf{NOT SUPPORTED} & 3 of 16 sets below 0.90: x flip 0.500 [0.388, 0.613], y flip 0.487, 10 Hz 0.263; base 0.963 & Impact direction from one IMU is the repository's own rule (\texttt{direction.py}); crash-pulse recorders exist (a 1995 validation study, seen in a listing only and not verified) & On simulated hardware-format logs the side rule failed under an axis-sign flip and at 10 Hz sampling and held 0.90 under tilts up to 30 degrees, noise x4 and bias 0.08 g. & supplement \\
H-S2 & The 2 g trigger gives crash recall of at least 0.90 and a non-crash false-alarm rate of at most 0.10 in every sensor variant. & SIMULATED\_HW\_FORMAT & no & 1 & 1 & \textbf{NOT SUPPORTED} & false-alarm rate of the 2 g trigger 0.640 [0.540, 0.740] in the base variant; recall ok in 15 of 16 sets (not at 10 Hz) & none needed (a check of a threshold on bench events) & On simulated bench events the 2 g trigger alone does not separate crash from non-crash (false-alarm rate about 0.6); this says nothing about vehicle crashes. & supplement \\
H-S3 & On link logs calibrated from the four edge runs, UA-SCT delivers at least as often as fixed-3 minus 0.02, within the duty allowance. & SIMULATED\_HW\_FORMAT & yes & 0.39 & 1 & \textbf{NOT SUPPORTED} & UA-SCT minus fixed-3 at 30, 34, 38, 42 dB: -0.0008, -0.0112 [-0.072, 0.050], 0.0008, 0.0000; largest airtime share 0.127 of the allowance & Tamang et al. 2023 (slotted alarms), Khan et al. 2019 (retry model), Sisinni et al. 2021 and 2022 (redundancy): baselines, not duplicated & On link logs simulated from stated models, no difference between UA-SCT and fixed-3 critical reach within 10 s was shown to be smaller than 0.02 (the 34 dB interval is wide); the severity table reached the base more often within 10 s at 34 dB. & supplement; negative register \\
H-S4 & The median trigger-to-ACK latency of first-try deliveries is at least 500 ms (a sub-500 ms statement fails for this window, alert and ACK). & SIMULATED\_HW\_FORMAT & yes & 1.26e-29 & 6.31e-29 & \textbf{SUPPORTED} & median 519 ms, p95 527 ms; lower bound 480.9 ms before compute & Carneiro et al. 2023 (LoRa crash response, record only); the sub-500 ms style claim is not attributed to a source & In the model, trigger-to-ACK latency is about 520 ms for a 250 ms post-trigger window, a 97-byte alert and a 26-byte ACK at SF7; each retry adds about 600 ms. This is a statement about the model, not the radio. & supplement (methods note) \\
H-S5 & No attack of the repository's suite and the simulated scenarios is accepted, and the reference verifier persists its counter. & SIMULATED\_HW\_FORMAT & yes & 3.65e-07 & 1.46e-06 & \textbf{SUPPORTED} & 0 of 289 attack attempts accepted (95\% upper bound 1.03\%); suite 58 cases, all as expected & Dong et al. 2024 and Wesemeyer et al. 2020 verify LoRaWAN; the alert suite is a different small protocol & The repository's attack suite and the nine simulated scenarios behave as designed in software; the ESP32 firmware was not tested (NOT\_RUN). & supplement (methods) \\
H-S6 & SW-420 confirmation lowers the non-crash false-alarm rate by at least 0.02 and loses at most 0.02 of the crash recall. & SIMULATED\_HW\_FORMAT & no & 2.25e-131 & 1.35e-130 & \textbf{NOT SUPPORTED} & false-alarm reduction 0.289 [0.266, 0.311], recall loss 0.043 [0.033, 0.053]; the sign of the criterion flips with the switch threshold range & a vibration-sensor accident paper of the project literature (08, SW-420 in three levels; \texttt{CLAIM\_EVIDENCE.md} N7); no benchmark found & In simulated bench events SW-420 confirmation removed 0.29 of the false alarms and cost 0.04 of the crash recall; the balance depends on the assumed switch sensitivity. & supplement; negative register \\
H-R1a & After a car reboot GTF accepts at least 99\% of valid alerts (clock drift up to 5 s, GPS loss up to 40\%). & SIMULATED & yes & 0.000745 & 0.00298 & \textbf{SUPPORTED} & smallest availability 0.9965 over 16 held-out cells; W = 10 s & ZigBee timestamp replay protection (IEEE TMC 2020, not opened); LoRaWAN frame counters; a LoRaWAN replay-attack study & In the simulation, a GPS-time freshness field accepted all valid alerts after a car reboot for clock drift up to 5 s with a 10 s window. & supplement \\
H-R1b & GTF replay acceptance is zero outside the window and at most 1\% inside it, at the held-out burst loads. & SIMULATED & no & 1 & 1 & \textbf{NOT SUPPORTED} & largest replay rate inside W 0.0325 (cells over 1\%: D2\_E2\_g0.2\_b300,D2\_E2\_g0.4\_b300,D2\_E4\_g0.2\_b300,D2\_E4\_g0.4\_b300,D5\_E2\_g0.2\_b300,D5\_E2\_g0.4\_b300,D5\_E4\_g0.2\_b300,D5\_E4\_g0.4\_b300); with a 512-entry ring 0.0000 & as R1a & The ring size chosen on the design grid (256) was too small for the held-out burst of 300 alerts per window (rate 1.7 to 3.0\%); a 512-entry ring met the bound. The ring must be sized above the number of alerts per window. & supplement; negative register \\
H-R1c & GTF cuts flash writes per year at least 10 times against persisting every alert counter. & ANALYTICAL & yes & 0 & 0 & \textbf{SUPPORTED} & 16.7x at 100 alerts and 12 boots per car and year (sensitivity table in \texttt{ar.r1.hR1c.sensitivity.*}) & none specific & Under stated rates, persisting no counter cuts flash writes about 17 times; the ratio is below 10 for fewer than 60 alerts per car and year at 12 boots. & supplement \\
H-R2a & A frame of at most 40 bytes with a 7 or 8 byte tag raises burst delivery by 5 points over the 97-byte frame. & SIMULATED & yes & 0.227 & 0.454 & \textbf{NOT SUPPORTED} & chosen on design: 24/8; 10 of 16 held-out cells gain at least 5 points; mean gain 0.075 (exploratory pooled 0.075 [0.068, 0.083]); the sign test over 16 cells gives p 0.227 & FC-HMAC (Zenodo record not read); the 4-byte LoRaWAN MIC; no study on tag length against burst delivery of single-packet alerts found in the sources checked & In simulated bursts a 24-byte frame delivered more critical alerts within 10 s than the 97-byte frame; the preregistered sign test over 16 cells did not reach significance, and the shorter frame drops fields (position) that were not designed here. & supplement; negative register (criterion met, test not significant) \\
H-R2b & The 97-byte frame with an 8-byte tag is dominated, in forgery effort, time on air and delivery, by a frame with a tag of at least 8 bytes. & SIMULATED & yes & 0 & 0 & \textbf{SUPPORTED} & dominating frames (tag >= 8 bytes): 24/8, 32/8, 32/12, 40/8, 40/12, 56/8, 56/12; point-estimate Pareto set (tag >= 7): 24/8, 32/12 & as R2a & In the simulation the 97-byte frame with an 8-byte tag is not on the Pareto front of forgery effort, time on air and burst delivery. & supplement \\
H-R3 & A provisional flash from the early-exit rule reaches the base at least 0.8 s earlier, with under 20\% retractions, under 40\% extra airtime and no delivery loss. & SIMULATED & yes & 0.0045 & 0.0135 & \textbf{SUPPORTED} & earlier by 1.87 s [1.87, 1.87] at loss 0; retractions 0.108 [0.096, 0.119]; extra airtime 0.380 [0.366, 0.395] at loss 0 & No crash-alert version of a provisional-then-confirm design found in the sources checked (not a claim of absence); early-exit classification of time series exists (Achenchabe et al., opened) & In simulation on de-duplicated VZCrash test events, a 14-byte provisional alert from the early-exit rule reached the base about 1.9 s before the full-window alert, 11\% of provisional alerts were retracted and the extra airtime was 38\%. Timing is the window length, not measured on the ESP32. The single-stage baseline waits for the full window; an early-exit single-stage alert would be as fast but has no cancel path. & main paper \\
H-R4 & With per-vehicle ACK keys a forged ACK from a compromised node stops only that node, at one extra KDF per new id and no extra bytes. & ANALYTICAL & yes & 0 & 0 & \textbf{SUPPORTED} & suppressed fraction 1.0 (shared key, event ACK) to 0.333 = 1/3 with one compromised vehicle, none of the honest vehicles; 0 extra bytes; one KDF per new id & Per-device keys are standard in LoRaWAN; the shared-key counterexample is from \texttt{pa\_model.py} & A bounded-model check found that a forged event-level ACK from a key holder stops every vehicle; per-vehicle keys with the ACK bound to the vehicle and counter limit the effect to the compromised vehicle. This is a design correction, not a new key scheme; a holder of the master key is outside the claim. & main paper \\
\bottomrule
\end{longtable}}

\subsection*{Candidate level}

{\scriptsize\raggedright
\begin{longtable}{p{0.237\linewidth}p{0.237\linewidth}p{0.237\linewidth}p{0.237\linewidth}}
\toprule
Candidate & Hypotheses & Verdicts & Placement \\
\midrule
\endhead
R1 reboot-resilient freshness (GTF) & H-R1a, H-R1b, H-R1c & SUPPORTED, NOT SUPPORTED, SUPPORTED & supplement; H-R1b to the negatives register (ring sizing) \\
R2 security-airtime co-design & H-R2a, H-R2b & NOT SUPPORTED, SUPPORTED & supplement; H-R2a to the negatives register \\
R3 two-stage alert & H-R3 & SUPPORTED & main paper \\
R4 per-vehicle ACK keys & H-R4 & SUPPORTED & main paper \\
\bottomrule
\end{longtable}}

Kill criteria: none fired. R1: drift of 2 s or less costs no availability and the time-stamped replay window is bounded by W. R2: a gain of at least 5 points exists and non-dominated frames exist. R3: all three clauses held. R4: the model found no other path that stops all vehicles from one compromised node. A kill criterion is not the same as a SUPPORTED verdict: H-R1b and H-R2a failed their preregistered criterion or test.

